\medskip

Les douanes s'intéressent aux importations de casques audio portant le logo d'une certaine marque. Les saisies des douanes permettent d'estimer que:

\setlength\parindent{10mm}
\begin{itemize}
\item[$\bullet~~$] 20\,\% des casques audio portant le logo de cette marque sont des contrefaçons ;
\item[$\bullet~~$] 2\,\% des casques non contrefaits présentent un défaut de conception ;\item[$\bullet~~$] 10\,\% des casques contrefaits présentent un défaut de conception.
\end{itemize}
\setlength\parindent{0mm}

L'agence des fraudes commande au hasard sur un site internet un casque affichant le logo de la marque. On considère les évènements suivants:

\setlength\parindent{10mm}
\begin{itemize}
\item[$\bullet~~$] $C$: \og le casque est contrefait \fg{} ;
\item[$\bullet~~$] $D$: \og le casque présente un défaut de conception \fg{} ;
\item[$\bullet~~$] $\overline{C}$ et $\overline{D}$ désignent respectivement les évènements contraires de $C$ et $D$.
\end{itemize}
\setlength\parindent{0mm}

Dans l'ensemble de l'exercice, les probabilités seront arrondies à $10^{-3}$ si nécessaire.

\bigskip

\textbf{Partie 1}

\medskip

\begin{enumerate}
\item Calculer $P(C \cap D)$. On pourra s'appuyer sur un arbre pondéré.
\item Démontrer que $P(D) = 0,036$.
\item Le casque a un défaut. Quelle est la probabilité qu'il soit contrefait ?
\end{enumerate}

\bigskip

\textbf{Partie 2}

\medskip

On commande $n$ casques portant le logo de cette marque. On assimile cette expérience à
un tirage aléatoire avec remise. On note $X$ la variable aléatoire qui donne le nombre de casques présentant un défaut de conception dans ce lot.

\medskip

\begin{enumerate}
\item Dans cette question, $n = 35$.
	\begin{enumerate}
		\item Justifier que $X$ suit une loi binomiale $\mathcal{B}(n,~p)$ où $n = 35$ et $p = 0,036$.
		\item Calculer la probabilité qu'il y ait parmi les casques commandés, exactement un casque présentant un défaut de conception.
		\item Calculer $P(X \leqslant 1)$.
	\end{enumerate}	
\item Dans cette question, $n$ n'est pas fixé.

Quel doit être le nombre minimal de casques à commander pour que la probabilité 
 qu'au moins un casque présente un défaut soit supérieur à $0,99$ ?
\end{enumerate}

\bigskip
